\begin{table}[t]
\centering
\caption{P3c M2: area\_total 분해 — 비조건부 (Ch4+Ch5) vs 조건부 (Ch4 단독)}
\label{tab:ch5_area_decomposition_ko}
\begin{tabular}{llrrr}
\toprule
Bandwidth & Year & Unconditional & Conditional (n=5) & Selection share \\
\midrule
T1 & 2018 & -149 & -138 & 7.1% \\
T1 & 2020 & -247 & -7 & 97.2% \\
T1 & 2021 & 44 & 131 & -197.3% \\
T1 & 2022 & -250 & -50 & 80.0% \\
T2 & 2018 & -42 & -100 & -136.2% \\
T2 & 2020 & 118 & 217 & -83.8% \\
T2 & 2021 & 259 & 253 & 2.4% \\
T2 & 2022 & 265 & 284 & -6.9% \\
T3 & 2018 & 46 & 58 & -26.5% \\
T3 & 2020 & 106 & 119 & -11.8% \\
T3 & 2021 & 313** & 271* & 13.4% \\
T3 & 2022 & 366*** & 350** & 4.5% \\
\bottomrule
\end{tabular}\\
\footnotesize\textit{비조건부 컬럼 = ch3\_area\_events (06\_channels.R Phase 2, 전체 패널). 조건부 컬럼 = n\_years==5 완전 농가만 (N=2,217). 선택 비중 = 1 - (조건부 / 비조건부). Channel 4(iii) 잔존 농가 내 확장 마진 = 조건부; Channel 5 선택 편의 = 비조건부 - 조건부 비율. * p$<$0.10, ** p$<$0.05, *** p$<$0.01.}
\end{table}

